\subsection{Exact neighboring-ray capacity}
\label{subsec:new-neighbor-ray-capacity}

The T3-like and Vd placements require the exact capacity of a V triangle on a
neighboring radial arm.  Normalize at $V_0$, let $A_a\in e_{5,0}$ and
$B_b\in e_{0,1}$ have reaches $a,b$, and let
\[
 X_c=(1-c)V_1\in r_1.
\]
Define $C_+(a,b)$ as the maximal $c$ for which
$V_0,A_a,B_b,X_c$ lie in a closed unit equilateral triangle.  Put
$C_-(a,b)=C_+(b,a)$.

For $0\le a\le1/2$, let $p(a)$ be the unique root in $[a,1]$ of
\[
 p^3-(a+2)p^2+2(a+1)p-1=0,
 \tag{6.24a}
\]
and define
\[
 \sigma(a)=1-p(a),
 \qquad
 \tau(a)=1-a-(p(a)-a)(1-p(a)).
\]

\begin{proposition}[Neighboring-ray formula]
\label{prop:new-neighbor-ray-formula}
\[
C_+(a,b)=
\begin{cases}
\text{undefined},&a+b>1,\\
1-b,&a+b\le1,\ a>1/2,\\
1-b,&0\le a\le1/2,\ 0\le b\le\sigma(a),\\
p(a),&0\le a\le1/2,\ \sigma(a)\le b\le\tau(a),\\
a+\dfrac12-\sqrt{(a+b)^2-\dfrac34},
 &0\le a\le1/2,\ \tau(a)<b\le1-a.
\end{cases}
\tag{6.24b}
\]
At $b=\tau(a)$ the plateau value $p(a)$ is selected.
\end{proposition}

\begin{proof}
Use the cone basis
\[
 E_-=V_5-V_0,
 \qquad E_+=V_1-V_0.
\]
In these coordinates
\[
 V_0=(0,0),\quad A=(a,0),\quad B=(0,b),\quad X=(c,1),
\]
and the metric is $u^2+v^2-uv$.

Let $R,S$ satisfy $R^2+S^2-RS=1$, and define the support coordinates
\[
 U(u,v)=(R-S)u+Sv,
\]
\[
 V(u,v)=Su-Rv,
 \qquad
 W(u,v)=Ru+(S-R)v.
\]
For this orientation, the least enclosing equilateral side is
\[
 \max_YW(Y)-\min_YU(Y)-\min_YV(Y),
 \tag{6.24c}
\]
where $Y$ ranges over $V_0,A,B,X$.  At a maximal feasible $c$, one support
side contains a hull edge.  The non-dominated possibilities are:
\[
\begin{array}{c|c}
\text{supported edges or contacts}&\text{candidate}\\
\hline
V_0A\text{ and }BX&c=1-b,\\
AX\text{ supported, }B\text{ inactive}&c=p(a),\\
AX\text{ supported, }B\text{ active}
 &c=a+\dfrac12-\sqrt{(a+b)^2-\dfrac34}.
\end{array}
\tag{6.24d}
\]
We verify the three patterns and their selectors.

For the first pattern take $R=S=1$.  The containing triangle is
\[
 U\ge0,
 \qquad V\ge-b,
 \qquad W\le c,
\]
with side $b+c$.  It contains $A$ exactly when $a\le c$.  Hence the unit
value is $c=1-b$, feasible precisely when $a+b\le1$.

For the second pattern put $d=c-a$.  The side through $A,X$ gives $R=Sd$,
and the unit condition gives $S^2(d^2-d+1)=1$.  With $B$ inactive, the
opposite support is $W(V_0)=0$ and $W(X)=-1$.  Eliminating $S$ gives
\[
(c-a)
\bigl(c^3-(a+2)c^2+2(a+1)c-1\bigr)=0.
\]
The nondegenerate root is $c=p(a)$.  For $0\le a\le1/2$ it is unique because
the cubic is strictly increasing on $[a,1]$.  The remaining containment
condition is
\[
 b\le1-p(a)(1+a-p(a))=\tau(a).
\]
This branch dominates $1-b$ exactly when $b\ge\sigma(a)$.

For the third pattern the active supports give
\[
 S=-\frac1{a+b},
 \qquad
 d^2-d+1=(a+b)^2.
\]
The selected root $0\le d\le1/2$ is
\[
 d=\frac{1-\sqrt{4(a+b)^2-3}}2,
\]
which gives the radical candidate in (6.24b).  The inactive support
inequalities reduce to
\[
0\le a\le\frac12,
\qquad
\tau(a)\le b\le1-a.
\]
On this interval the radical candidate dominates $1-b$ because, with
$s=a+b\le1$,
\[
 s-\frac12\ge\sqrt{s^2-\frac34}.
\]

The three selected intervals are ordered since
\[
 \tau(a)-\sigma(a)=(p(a)-a)p(a)\ge0,
 \qquad
 \tau(a)\le1-a.
\]
If $a+b>1$, none of the three support patterns is feasible, including
$c=0$, so the capacity is undefined.  Reflection gives $C_-$.
\end{proof}

A consequence used repeatedly below is
\[
 \boxed{C_+(p,q),C_-(p,q)\le1-\min\{p,q\}}
 \qquad(p+q\le1).
 \tag{6.24e}
\]
Indeed this is immediate on the linear branch; on the plateau branch the
cubic gives $p(a)\le1-a$, and the radical branch is no larger than its
selected plateau at the transition.
